\chapter*{安装}
\addcontentsline{toc}{chapter}{安装}

运行原书代码需要 Python~3、PyTorch 与辅助包 \code{d2l}。
记依赖集合为 $\mathcal{P}$。把解释器版本与 $\mathcal{P}$ 配成同一隔离环境后，
笔记本里的张量运算才能按后文章节的公式求值；配错版本则同一符号可能对应不同实现。

原书推荐用 conda 创建名为 \code{d2l} 的环境，再安装 CPU 或 GPU 版 PyTorch。
有 NVIDIA GPU 且已配置 CUDA 时，同一前向
\[
\hat{y}=f_{\bm{\theta}}(\bm{x})
\]
与同一反向 $\nabla_{\bm{\theta}}L$ 可在加速器上求值：算出的 $\hat{y}$、$L$ 与各梯度分量
在数学含义上与 CPU 相同，只是求值更快。前几章用 CPU 即可完成。

每次运行前激活该环境。激活成功表示后续 \code{import torch} 读到的是这一套包，
而不是系统里另一份；它本身不是一个数值，而是“公式将在哪一套实现上被计算”的选择。
